\chapter{Configuration Philosophy}

Configuration is cognitive architecture in code form.

\section{Permission Levels}

\begin{table}[H]
\centering
\begin{tabular}{lll}
\toprule
\textbf{Level} & \textbf{When to Use} & \textbf{DCF Implication} \\
\midrule
Ask always & Learning, sensitive code & Maximum DCF engagement \\
Allow with approval & Normal development & Checkpoint-based DCF \\
Allow automatically & Routine, trusted operations & Trust but verify results \\
\bottomrule
\end{tabular}
\caption{Permission Levels and DCF}
\end{table}

\textbf{The Configuration DCF Principle:} Your configuration should match your cognitive engagement level. More autonomy = more need for post-hoc review. Less autonomy = more in-flow checkpoints.

\section{Hooks: Automated DCF Triggers}\index{Hooks}\label{sec:hooks}

Hooks are shell commands that execute automatically before or after Claude Code tool calls. They enable \textbf{automated DCF checkpoints}.

\subsection{Hook Types}

\begin{table}[H]
\centering
\begin{tabular}{lll}
\toprule
\textbf{Hook Type} & \textbf{Trigger} & \textbf{DCF Use Case} \\
\midrule
PreToolCall & Before any tool runs & Validation, logging \\
PostToolCall & After tool completes & Auto-testing, verification \\
Notification & On specific events & Alerting for review \\
\bottomrule
\end{tabular}
\caption{Hook Types and DCF Applications}
\end{table}

\subsection{DCF-Informed Hook Patterns}

\begin{lstlisting}
# .claude/hooks/post-edit.sh
# Run tests after any file edit - automated verification
npm test --related $CHANGED_FILES

# .claude/hooks/pre-commit.sh
# Require human review before commits
echo "DCF Checkpoint: Review changes before commit"
git diff --staged
\end{lstlisting}

\textbf{The Hook Philosophy:} Hooks externalize your DCF principles into automated enforcement. If you always want tests run after edits, don't rely on memory---encode it in a hook.

\section{Settings Configuration}\index{Settings}

The \texttt{.claude/settings.json} file defines project-level behavior. Key settings from a DCF perspective:

\begin{lstlisting}
{
  "permissions": {
    "allow": ["Read", "Glob", "Grep"],
    "ask": ["Edit", "Write", "Bash"],
    "deny": ["dangerous-commands"]
  },
  "planMode": {
    "requireApproval": true,
    "maxTurns": 50
  }
}
\end{lstlisting}

\subsection{Permission Strategy}

\begin{itemize}
    \item \textbf{allow}: Tools that don't modify state---read operations are generally safe
    \item \textbf{ask}: Tools that change files or run commands---checkpoint opportunities
    \item \textbf{deny}: Operations you never want automated
\end{itemize}

\textbf{DCF Principle:} Configure permissions to create natural checkpoints at high-stakes moments while allowing flow for low-stakes operations.

\section{Model Selection as DCF Decision}\index{Model Selection}

Claude Code supports multiple models with different capabilities and costs. Model selection is itself a DCF decision.

\begin{table}[H]
\centering
\begin{tabular}{llll}
\toprule
\textbf{Model} & \textbf{Strength} & \textbf{Cost} & \textbf{When to Use} \\
\midrule
Haiku & Speed, simple tasks & Lowest & File searches, quick questions \\
Sonnet & Balance & Medium & Most development work \\
Opus & Deep reasoning & Highest & Architecture, complex judgment \\
\bottomrule
\end{tabular}
\caption{Model Selection Guide}
\end{table}

\subsection{DCF Model Selection Heuristics}

\begin{itemize}
    \item \textbf{Use Haiku} for exploration, searches, and tasks where you'll heavily review output anyway
    \item \textbf{Use Sonnet} for standard development where balance matters
    \item \textbf{Use Opus} when the decision is hard to reverse, architecturally significant, or requires nuanced judgment
\end{itemize}

The principle: \textbf{match model capability to decision stakes}. Routine operations don't need your most powerful model; critical decisions deserve it.

\section{Session Lifecycle}\index{Session Lifecycle}\label{sec:session-lifecycle}

Knowing when to start a fresh session vs. continue an existing one is a form of cognitive hygiene.

\subsection{When to Start Fresh}

\begin{itemize}
    \item Context has drifted significantly from original goal
    \item Accumulated errors or misunderstandings are compounding
    \item You're starting a genuinely new task (not a continuation)
    \item Session history is no longer helpful (noise > signal)
\end{itemize}

\subsection{When to Continue}

\begin{itemize}
    \item Task is directly related to prior work
    \item Context from prior exchanges is valuable
    \item You haven't hit compaction yet and flow is good
    \item The AI has learned project-specific patterns you'd lose
\end{itemize}

\subsection{Session Hygiene Practices}

\begin{enumerate}
    \item \textbf{Name your sessions mentally}---``This is the API refactor session''
    \item \textbf{Use /compact strategically}---when context is large but still valuable
    \item \textbf{Capture before switching}---use \texttt{/dcf compact} before major context shifts
    \item \textbf{Start fresh for fresh tasks}---don't let inertia keep you in a stale session
\end{enumerate}

\subsection{The 50-60\% Context Guideline}

A practical heuristic: \textbf{effective context capacity is 50-60\% of the maximum, not 100\%}.

This matters because:
\begin{itemize}
    \item \textbf{Tool calls consume context}: Every file read, command executed, and search result adds to context consumption---often invisibly.
    \item \textbf{Performance degrades before limits}: AI quality can degrade well before you hit hard context limits.
    \item \textbf{Headroom enables recovery}: If you're at 90\% and need to paste an error message, you have no room to work.
\end{itemize}

Practical implications:
\begin{itemize}
    \item \textbf{Monitor context}: Use \texttt{/context} periodically to check consumption.
    \item \textbf{Plan for compaction early}: When you hit 50-60\% on a complex task, start thinking about what to capture.
    \item \textbf{Don't start complex tasks mid-session}: If you're already at 40\%, starting a major new task risks hitting limits at an awkward point.
    \item \textbf{Summarize rather than paste}: A 500-line document becomes noise; a 10-line summary with ``full doc at path/to/spec.md'' preserves headroom.
\end{itemize}

The 50-60\% guideline isn't a rigid rule, but it prevents the surprise of hitting context limits during critical work.

\section{IDE Integration}\index{IDE Integration}

Claude Code integrates with development environments, extending DCF into your editing workflow.

\subsection{VS Code Extension}

The Claude Code VS Code extension provides:
\begin{itemize}
    \item Inline chat within the editor
    \item File context automatically included
    \item Code selection as conversation input
    \item Terminal integration for command execution
\end{itemize}

\subsection{DCF in IDE Context}

IDE integration changes the DCF dynamic:

\begin{table}[H]
\centering
\begin{tabular}{lll}
\toprule
\textbf{Aspect} & \textbf{CLI} & \textbf{IDE} \\
\midrule
Context & You provide files explicitly & Files visible automatically \\
Review & Terminal output, then check files & Changes visible in diff view \\
Flow & Conversation-centric & Code-centric with chat \\
Checkpoints & Before commands execute & Before accepting changes \\
\bottomrule
\end{tabular}
\caption{CLI vs. IDE Integration}
\end{table}

\subsection{IDE DCF Practices}

\begin{itemize}
    \item \textbf{Use diff view as your mirror}---review proposed changes before accepting
    \item \textbf{Select relevant code}---don't rely on the AI to find the right context
    \item \textbf{Checkpoint at accept}---the ``Accept'' button is your approval point
    \item \textbf{Keep terminal visible}---command execution should still be reviewed
\end{itemize}

\textbf{The IDE Principle}: IDE integration makes some context automatic but doesn't remove the need for engagement. The convenience is in context gathering; the judgment remains yours.

